% Part: lambda-calculus
% Chapter: introduction
% Section: composition

\documentclass[../../../include/open-logic-section]{subfiles}

\begin{document}

\olfileid[tr]{lam}{int}{com}
\olsection{\usetoken{S}{lambda definable} Fonksiyonlar Bileşim Altında Kapalıdır}

\begin{lem}
!!{lambda definable} fonksiyonlar bileşim altında kapalıdır.
\end{lem}

\begin{proof}
$f$'nin $h$, $g_0$, \dots, $g_{k-1}$'den bileşimle tanımlandığını
varsayalım. $h$, $g_0$, \dots, $g_{k-1}$'in sırasıyla $H$, $G_0$, \dots,
$G_{k-1}$ tarafından !!{lambda defined} olduğu varsayımı altında $f$'yi
!!{lambda define}s bir $F$ terimi bulmalıyız. Fakat $F$'yi doğrudan şöyle
tanımlayabiliriz:
\[
F(x_0, \dots, x_{l-1}) = H(G_0(x_0, \dots, x_{l-1}),
\dots, G_{k-1}(x_0, \dots, x_{l-1})).
\]
Başka bir deyişle lambda kalkülüsünün dili bileşimi temsil etmeye çok
uygundur.
\end{proof}

\end{document}
